% Source: https://artofproblemsolving.com/community/c3h1883620_probability_and_geo
% Stable identity: aops-topic-1883620-q01
Place the square at $[0,2]\times[0,2]$ and let the point be $P=(x,y)$, uniformly distributed in the square. Each triangle has a side of the square (length $2$) as base, and its height is the distance from $P$ to that side, so the four areas are
\[
y,\quad 2-y,\quad x,\quad 2-x
\]
(bottom, top, left and right triangles).

\begin{asy}
usepackage("amsmath");
size(5cm);
defaultpen(fontsize(10pt));
pair A=(0,0),B=(2,0),C=(2,2),D=(0,2),P=(1.25,0.8);
real t=2/3;
pen sh=rgb(0.85,0.92,1);
fill(A--(t,0)--(0,t)--cycle,sh);fill(B--(2-t,0)--(2,t)--cycle,sh);
fill(C--(2,2-t)--(2-t,2)--cycle,sh);fill(D--(t,2)--(0,2-t)--cycle,sh);
draw((t,0)--(0,t),blue);draw((2-t,0)--(2,t),blue);draw((2,2-t)--(2-t,2),blue);draw((t,2)--(0,2-t),blue);
draw(A--B--C--D--cycle,linewidth(1));
draw(P--A);draw(P--B);draw(P--C);draw(P--D);
draw(P--(1.25,0),dashed);draw(P--(0,0.8),dashed);
dot(P);label("$P$",P,N);
label("$x$",(0.62,0.8),N,fontsize(9pt));label("$y$",(1.25,0.4),E,fontsize(9pt));
label("$2$",(0,1),W);
label("$\dfrac{2}{3}$",(t/2,0),S,fontsize(8pt));
label("\mbox{}",(min(currentpicture,user=true)+max(currentpicture,user=true))/2,black); // keeps the text after the figure black when LaTeX typesets the labels inline
\end{asy}

Two opposite triangles always have total area $2$. Two adjacent triangles have total area $x+y$, $(2-x)+y$, $(2-x)+(2-y)$ or $x+(2-y)$, i.e.\ the sum of the distances from $P$ to two adjacent sides. So the condition holds exactly when $P$ lies in one of the four corner triangles cut off by the lines $x+y=\dfrac{2}{3}$, etc. (shaded). Each is a right isosceles triangle with legs $\dfrac{2}{3}$ and area $\dfrac{2}{9}$, and they do not overlap. Hence the probability is
\[
\dfrac{4\cdot\dfrac{2}{9}}{4}=\boxed{\dfrac{2}{9}}.
\]
